%%
% BIThesis 本科毕业设计论文模板 —— 使用 XeLaTeX 编译 The BIThesis Template for Undergraduate Thesis
% This file has no copyright assigned and is placed in the Public Domain.
%%

\begin{conclusion}
  % 结论部分尽量不使用 \subsection 二级标题，只使用 \section 一级标题

  本文围绕自重组轮足机器人链式系统协同控制问题，开展了单体建模、平衡控制和链式协同运动建模研究。首先，针对单体轮腿机器人结构，建立了腿部运动学模型和虚拟模型控制映射关系，推导了腿部虚拟支撑力、摆杆力矩与关节驱动力矩之间的等效关系。其次，针对尾部辅助机构离地和触地两类典型工况，建立了单体动力学模型，并在平衡点附近线性化得到状态空间表达。在此基础上，设计了 LQR 状态反馈控制器，并通过参数化增益拟合方法提高控制器对腿长变化和构型变化的适应能力。

  针对链式系统协同运动问题，本文进一步分析了相邻单体之间的铰接几何约束，建立了由头车运动状态向后续单体逐级传播的运动学模型。基于该模型，构建了分层协同控制思路：上层根据链式几何约束生成各单体期望运动序列，下层利用单体平衡控制器实现稳定跟踪。该方法将链式协同运动问题与单体姿态稳定问题相分离，有助于降低控制器设计复杂度，并为后续引入 DMPC 约束优化提供了基础。

  后续研究将重点从三个方面继续完善：一是开展系统化实机实验，形成可量化的速度误差、姿态误差、稳定时间和控制输入对比结果；二是完善 DMPC 协同层，使其能够在预测时域内显式处理铰接角、速度和输入约束；三是进一步研究 LQR-ADRC 复合控制，通过扰动观测和补偿提升单体在模型误差、冲击载荷和复杂地形下的鲁棒性。
\end{conclusion}
